\section{The question this report answers}
\label{sec:question}

The previous report evaluated the multilayer convolutional Tsetlin machine --- convolutional
layers that keep their spatial clause-activation maps instead of OR-pooling them away,
trained greedily and frozen --- and split cleanly into a positive and a negative result.

\begin{ttkey}
\lead{The architecture was vindicated.} On a task whose one-layer ceiling is known
analytically, a two-layer stack given a hand-built first layer cleared it
($\synXorMctmOracleSixFour{}\%$ against $75\%$ at a matched channel count); on a relational
task the variant that keeps the spatial map reached $\synNearMctmTask{}\%$ while the
global-pooling stack was pinned at chance whatever features it was given.

\lead{The training scheme was not.} On CIFAR-10 the stack reached \accMctm{}\% against
\accSingleMatched{}\% for a single layer of the same total size, and fell below the
\accSingleLOne{}\% its own frozen first layer achieved alone. \emph{Random} first-layer
clauses beat trained ones, \accMctmRandom{}\% against \accMctm{}\%.
\end{ttkey}

The diagnosis was a coupling. Training layer~1 to classify rewards \emph{specific} clauses ---
many included literals, rare firing --- while layer~2 needs \emph{general} ones whose
activations are common enough to be conjoined. Layer~1 fires on a fraction of a percent of
patch positions, and layer~2, facing an input where positive literals are almost never
satisfied and negations almost always are, becomes \negMctm{}\% negated literals and learns
that things are absent.

That report's \S7.3 named three ways out, in increasing order of how much they change the proposal.
This report implements all three and measures them.

\begin{itemize}[leftmargin=1.2em]
\item[\ideaA] \lead{Give layer 1 a different objective.} Nothing forces a feature layer to be
  trained by classification. A reconstruction objective --- clauses that predict held-out
  features of a patch from the rest --- preserves information rather than discarding whatever
  did not help the local classifier, and has no reason to drive clauses towards extreme
  specificity.
\item[\ideaB] \lead{Calibrate density explicitly.} Rather than hoping the specificity
  parameter lands somewhere usable, fix a target per-channel firing rate and adjust each
  clause towards it --- a Boolean analogue of normalisation.
\item[\ideaC] \lead{Propagate credit instead of freezing.} Clause inclusion is a usable credit
  path: if a layer-2 clause receives Type~I feedback and includes the literal ``channel $k$
  active at offset $(dy,dx)$'', then layer-1 clause $k$ can receive Type~I feedback at that
  position. A discrete, rule-based analogue of backpropagation that requires no gradients.
\end{itemize}

\subsection{What this report adds}

Each idea is implemented, measured on CIFAR-10 against the same baselines at the same clause
budget, and measured again on the two constructed tasks of the first report, where the
per-arm ceiling is known in advance and the oracle says how much headroom a better training
scheme actually has.

Four things were not obvious before doing it, and they are the substance of what follows.

The unsupervised objective does exactly what it was predicted to do to the \emph{statistics}
of the representation, and on CIFAR-10 that is worth six points and no more. The density
controller reaches any target rate asked of it, and produces a monotone dose-response curve
that turns the first report's density argument into a measurement. Credit propagation, written
out honestly, turns out to have no forgetting term, then no bandwidth, and finally a drift
problem that is more interesting than either.

And the fourth is the one that matters. Run over an \emph{untrained} first layer --- random
clauses, never shown a label --- the same density controller takes CIFAR-10 to
\randBestAcc\,$\pm$\,\randBestSd{}\%, past the \accSingleMatched{}\% of a single layer at the
same clause budget and the \accSingleRf{}\% of a single layer given the stack's receptive
field. That is the first configuration in either report for which the second layer is worth
what it costs. On the constructed relational task, where a single layer is at chance
(\synNearSingleTask{}\%) by construction, repairing the first layer's density by either route
solves the task outright: \synNearAutoAuto{}\% with the autoencoder layer and
\synNearCalibTask{}\% with the controller, against \synNearMctmTask{}\% for greedy training
and \synNearMctmOracleSixFour{}\% for the \emph{hand-built oracle} of the first report, whose
two exact shape detectors turn out to be too sparse to conjoin.

So the first report's diagnosis was right and its repairs work. What does not survive is the
part of the proposal nobody was questioning: that layer~1 should be trained.
